\section{Introduction}
\label{sec:introduction}

General-purpose vision-language-action (VLA) models can learn a wide range of diverse manipulation skills from data. Yet they often struggle in the last millimeter of execution: motions can be slow, successful completion may require pauses and retries, and small errors at critical stages of precise tasks can compound into failure. A natural way to address this challenge is to fine-tune VLAs with reinforcement learning (RL). By practicing on the target task, RL can improve precisely the stages of the task that are most critical for success, which are often the stages most sensitive to small errors and the hardest to cover reliably with demonstrations alone. But real-world robotics operates under a tight budget: every episode takes time, every failure consumes effort and wear, and meaningful adaptation often has to happen within a few hours of practice.

Sample-efficient fine-tuning of VLAs, however, presents major challenges. On one hand, conventional methods for RL training of foundation models~\citep{Li2025SimpleVLA_RL,li2025grrlgoingdexterousprecise,pi06vla} rely on large-scale data and can be inefficient for rapid online adaptation. On the other hand, data-efficient real-world RL methods~\citep{luo2024precise,rl100} typically train much smaller models, which can improve within hours but sacrifice the generalization capabilities of the VLA. The central question, then, is how to leverage the VLA’s generalization while achieving the speed and sample efficiency of lightweight online RL.
%%SL.3.16: This seems like a false dichotomy. We could also simply delete this paragraph, but if you want to keep it, I would shorten it and avoid the "fundamental tradeoff" (it doesn't seem like a fundamental tradeoff).
%%KK.3.16: replaced fundamental tradeoff with "challenges". This was rewriting

We introduce a practical recipe that bootstraps fast online reinforcement learning with representations obtained from a pretrained VLA policy. Our key idea is to adapt the VLA so that it exposes a compact interface that can be used for sample-efficient online RL. We achieve this by training the VLA to expose an \emph{RL token}, a compressed representation that makes task-relevant pretrained knowledge accessible to a lightweight online RL policy. Running RL with this RL token (\MethodName{}) creates a simple division of labor: the frozen VLA provides broad perceptual understanding and action recommendations, while the lightweight actor and critic adapt the policy to succeed at the hardest parts of a task online. To make this practical in the sample-efficient real-world regime, our method uses a sample-efficient online RL algorithm to train the small actor and critic networks that use the RL token representation, with an additional regularizer to anchor the actor to the VLA action, so that online RL refines promising behaviors rather than learning from scratch.

We evaluate \MethodName{} on four challenging robot manipulation tasks that can require millimeter or sub-millimeter precision: screw installation, zip tie fastening, Ethernet, and charger insertion.
Across these tasks, \MethodName{} improves both success rate and execution speed within a few hours of online training. The largest gains appear in the critical phases of the tasks, which require high precision and determine task success, where \MethodName{} speeds up execution by up to 3$\times$ and substantially improves success rates, for example from 20\% to 65\% for a challenging screw-insertion task. On one of the most dexterous parts of our tasks, policies trained with our method can surpass expert teleoperation speed while maintaining reliability. These results suggest that combining a VLA model with lightweight online RL offers a practical path to high-performance manipulation without extensive task-specific engineering.
